\documentclass[11pt,a4paper]{article}
\usepackage[T1]{fontenc}
\usepackage{lmodern}
\usepackage[margin=27mm]{geometry}
\usepackage{amsmath,amssymb,amsthm,mathtools,booktabs,microtype}
\usepackage[hidelinks]{hyperref}
\hypersetup{pdftitle={Asymptotic enumeration of tied football seasons},pdfsubject={A380592 leading term and fixed-order asymptotic expansion}}
\setcounter{tocdepth}{1}
\usepackage{enumitem}
\setlist{itemsep=3pt,topsep=5pt}
\setlength{\parindent}{1.2em}
\setlength{\parskip}{2pt}
\setlength{\emergencystretch}{2em}
\allowdisplaybreaks[2]
\numberwithin{equation}{section}
\newtheorem{theorem}{Theorem}[section]
\newtheorem{lemma}[theorem]{Lemma}
\newtheorem{proposition}[theorem]{Proposition}
\theoremstyle{remark}
\newtheorem{remark}[theorem]{Remark}
\newcommand{\E}{\mathbb E}
\newcommand{\Prob}{\mathbb P}
\newcommand{\R}{\mathbb R}
\newcommand{\Z}{\mathbb Z}
\newcommand{\1}{\mathbf 1}
\newcommand{\ii}{\mathrm i}
\newcommand{\Var}{\operatorname{Var}}
\newcommand{\Cov}{\operatorname{Cov}}
\newcommand{\diag}{\operatorname{diag}}
\newcommand{\poly}{\operatorname{poly}}
\newcommand{\ceil}[1]{\left\lceil #1\right\rceil}
\newenvironment{writenote}[1][]{\par\smallskip\begingroup\small
 \noindent\textbf{[write]}\if\relax\detokenize{#1}\relax\else~\textbf{#1.}\fi\ }%
 {\par\endgroup\smallskip}
\title{Asymptotic enumeration of tied football seasons}
\author{}
\date{2 October 2026}
\begin{document}
\maketitle
\vspace{-1.5em}
\begin{abstract}
For a double round-robin league with labeled teams, distinguished matches,
and the usual $3/1/0$ scoring rule, let $a(n)$ count seasons in which every
team has the same score. We prove a leading equivalent for OEIS A380592,
including its nonconstant three-residue theta factor. The proof uses
mean-matching exponential tilts and a sharp global Fourier localization:
a negative quartic gives an exact local Gaussian majorant, and a separate
outlier argument controls its complement. The limiting common score is a
discrete Gaussian whose center differs from its mean. We then verify the
hypotheses of Isaev's complex cumulant theorem to obtain an expansion to
every fixed inverse-power order, give a finite exact coefficient algorithm,
compute the first correction, and formulate residue-aware inverse
threshold bounds. The supplied replay checks the finite algebra; the
analytic estimates and their scope are proved in the text.
\end{abstract}
\tableofcontents
\newpage

\section{Model and main results}\label{tfs:sec:model}
There are $n$ labeled teams. Every unordered pair plays two distinguished
matches. A match between $i$ and $j$ has score outcomes $(3,0)$, $(1,1)$,
and $(0,3)$. Put $M=n(n-1)$, the number of matches. Then
\begin{equation}\label{tfs:eq:exact}
a(n)=\sum_{k=2(n-1)}^{3(n-1)}
 [x_1^k\cdots x_n^k]\prod_{i<j}(x_i^3+x_ix_j+x_j^3)^2.
\end{equation}
The endpoints follow by summing all scores: each match contributes either
$2$ or $3$. This counts neither unlabeled leagues nor score tables alone.
If outcomes are independent and uniform and $X_i$ denotes team $i$'s score,
then $a(n)=3^M\Prob(X_1=\cdots=X_n)$.

The exact enumeration paper of Jehn, Habermann and Lavrov
\cite{football} studies this model and gives values through $n=8$.
The values in OEIS A380592 \cite{oeis} are
\[
1,\ 3,\ 27,\ 1083,\ 296081,\ 696779523,\
16503494334993,\ 3439079361325736243.
\]
Define the positive, real, $1$-periodic function
\begin{align}
\Theta(c)&=\sum_{j\in\Z}e^{-4\pi^2j^2/9}e^{2\pi\ii jc},\label{tfs:eq:theta}\\
&=\frac{3}{2\sqrt\pi}\sum_{k\in\Z}e^{-9(k-c)^2/4},\label{tfs:eq:poisson}
\end{align}
where the second equality is Poisson summation. Write
\begin{equation}\label{tfs:eq:carrier}
\mu_n=\frac{8(n-1)}3,\qquad c_n=\mu_n-\frac{13}{84},\qquad
B_n=3^M n^{-1/2}\left(\frac9{56\pi n}\right)^{(n-1)/2}
 e^{-9829/65856}.
\end{equation}

\begin{theorem}[Leading equivalent]\label{tfs:thm:leading}
As $n\to\infty$,
\begin{equation}\label{tfs:eq:main}
a(n)\sim B_n\Theta(c_n).
\end{equation}
Equivalently, retaining the exact transverse covariance,
\begin{equation}\label{tfs:eq:exactcarrier}
a(n)\sim 3^M(n-1)^{-1/2}
 \left(\frac9{2\pi(28n-2)}\right)^{(n-1)/2}
 e^{-12181/65856}\Theta(c_n).
\end{equation}
\end{theorem}
Because $c_{n+3}-c_n=8$, the multiplier depends only on $n\bmod 3$.
It cannot be replaced by $1$ in an asymptotic equivalent.

\begin{theorem}[Every fixed order]\label{tfs:thm:allorders}
For every fixed integer $R\geq0$, there are explicitly computable
$1$-periodic functions $H_0,\ldots,H_R$ such that
\begin{equation}\label{tfs:eq:allorders}
a(n)=B_n\left(\sum_{j=0}^R\frac{H_j(c_n)}{n^j}
                   +O_R(n^{-R-1})\right),\qquad H_0=\Theta.
\end{equation}
Each $H_j$ is a finite rational linear combination of derivatives of
$\Theta$. In particular, set
\begin{equation}\label{tfs:eq:Q}
Q(\delta)=\frac{22180735}{45177216}
 +\frac{98841}{153664}\delta-\frac{5283}{3136}\delta^2
 -\frac98\delta^3.
\end{equation}
For $r\in\{0,1,2\}$ and any $n\equiv r\pmod3$, define
\begin{equation}\label{tfs:eq:Rr}
R_r=\frac{\displaystyle\sum_{k\in\Z}e^{-9(k-c_n)^2/4}Q(k-\mu_n)}
           {\displaystyle\sum_{k\in\Z}e^{-9(k-c_n)^2/4}}.
\end{equation}
Then
\begin{equation}\label{tfs:eq:first}
a(n)=B_n\Theta(c_n)\left(1+\frac{R_{n\bmod3}}n+O(n^{-2})\right).
\end{equation}
\end{theorem}
The defining sums are exact. The following rounded decimals are numerical
diagnostics, not interval-certified enclosures.
\begin{center}
\begin{tabular}{@{}ccrr@{}}
\toprule
$n\bmod3$ & $c_n\pmod1$ & $\Theta(c_n)$ & $R_{n\bmod3}$\\
\midrule
$0$ & $5/28$ & $1.01079874779311$ & $0.04288551652736$\\
$1$ & $-13/84$ & $1.01402025457159$ & $0.19984197501234$\\
$2$ & $43/84$ & $0.97518099763530$ & $0.04601392878474$\\
\bottomrule
\end{tabular}
\end{center}
Sections~\ref{tfs:sec:tilt}--\ref{tfs:sec:summation} give an elementary proof of
Theorem~\ref{tfs:thm:leading}. Sections~\ref{tfs:sec:allorders}--\ref{tfs:sec:first}
prove Theorem~\ref{tfs:thm:allorders}, invoking a stated complex-cumulant theorem
with verified hypotheses. Section~\ref{tfs:sec:inverse} treats inversion.

\begin{writenote}[Added 2 October 2026, batch~77]
This report is a single manuscript, printed in full: manuscript~67 of
batch~77 of ProveIt's incoming-reports intake, delivered as the archive
\texttt{tied-\allowbreak football-\allowbreak report.zip} (wrapper directory
\texttt{tied-\allowbreak football-\allowbreak release/}, main file
\texttt{report/\allowbreak tied-\allowbreak football-\allowbreak asymptotics.tex},
now \texttt{article.tex}) in commit \texttt{096ee7b87}, and placed in
commit \texttt{d0e6008d9} in the research-report collection, among the
OEIS-sequence asymptotics reports. The manuscript names no ProveIt commit
and describes no search of the repository, so it has no pin; its author
line is empty. The bounded search of Section~\ref{tfs:sec:prior} is the
manuscript's own literature search; the intake's untruncated search of the
repository (the research-report collection, the transseries volumes and
\texttt{Oeis/}) found no other treatment of A380592 or of tied seasons.
The nearest batch-77 neighbour, the collection report
\texttt{a000571-\allowbreak tournament-\allowbreak score-\allowbreak sequences}
(batch-77 manuscript~68), counts score sequences of tournaments from their
exact generating function; it is a different model and does not use the
tilted Fourier localization of this report, and neither report uses the
other. Nothing in this report is formalized in Lean
or Rocq, and placement in the collection confers no formal status on any
statement here. The writing step prefixed every label with \texttt{tfs:},
labelled Section~\ref{tfs:sec:model}, added these \textsf{[write]} notes
and two references \cite{tfs-tai,tfs-cti}, and set the bibliography
ragged-right; no statement, proof or number of the manuscript was altered.

\emph{Notation.} Each symbol is printed as delivered; several letters carry
two meanings in different sections:
\begin{center}\footnotesize
\begin{tabular}{@{}p{0.15\linewidth}p{0.79\linewidth}@{}}
\toprule
symbol & meanings (and what each is \emph{not})\\
\midrule
$Q$, $Q_0$ & the first-correction polynomial $Q(\delta)$ of \eqref{tfs:eq:Q}; $Q_0$ is a small closed interval of tilt parameters $q$ about $1/3$ (Section~\ref{tfs:sec:localize}), not a value of $Q$.\\
$R$, $R_r$ & a truncation order in Theorem~\ref{tfs:thm:allorders} and Section~\ref{tfs:sec:inverse}; $R_r$ ($r=n\bmod3$) is the first-correction constant of \eqref{tfs:eq:Rr}.\\
$B_n$, $B$, $b_3$, $B_4$ & the carrier $B_n$ of \eqref{tfs:eq:carrier}; the bad-coordinate set $B$ of Section~\ref{tfs:sec:localize}; $b_3,B_4$ cubic and quartic coefficients in Section~\ref{tfs:sec:first}.\\
$b$, $d$, $e_r$ & in Section~\ref{tfs:sec:inverse}, $b=\log3$, $d$ and $e_r$ of \eqref{tfs:eq:de}. Elsewhere $d=(x-y)/2$ (Section~\ref{tfs:sec:localize}), $d(u)$ the torus distance, $d_3$ a cubic coefficient, and $D(q\Vert1/3)$ the Bernoulli relative entropy; $D$ is also a Taylor degree in Section~\ref{tfs:sec:allorders}.\\
$c$, $c_n$ & $c_n=\mu_n-13/84$ is the Gaussian centre of the common score; $c=na$ is the scaled common Fourier variable in Section~\ref{tfs:sec:local}; $c,c',c''$ are also positive constants. The centre $c_n$ is not the conditional mean (Proposition~\ref{tfs:prop:conditional}).\\
$K$, $K_0$ & $K_0=e^{-7865/32928}$ is the local correction \eqref{tfs:eq:K}; $K=\log n$ is the Gaussian truncation level in Section~\ref{tfs:sec:allorders}; $k$ is the common score.\\
$\Theta$, $\Theta_r$, $\theta$ & the theta factor \eqref{tfs:eq:theta} and its residue values $\Theta_r=\Theta(c_n)$; $\theta\in[-\pi,\pi]^n$ is the Fourier variable.\\
$\rho_n$, $\lambda_n$, $\alpha$, $\beta$ & the covariance eigenvalues \eqref{tfs:eq:eigen} and $\alpha(q)=2v(q)$, $\beta(q)=q(1-q)$; $v(q)$ differs from the covariance entry $v$ of \eqref{tfs:eq:coventries}, as $u$ there differs from the one-match variable $U$.\\
$M$, $M_j$ & $M=n(n-1)$ is the number of matches; $M_j=\E f_D^j$ are moments in \eqref{tfs:eq:cumrec}.\\
$L$, $L_r$ & $L=\log y$ is the inverse target in Section~\ref{tfs:sec:inverse}; $L_r$ are the degree-$r$ logarithmic polynomials of Sections~\ref{tfs:sec:allorders}--\ref{tfs:sec:first}.\\
$H_j$, $H_3$ & the periodic coefficient functions of Theorem~\ref{tfs:thm:allorders}; $H_3(z)=z^3-3z$ is a Hermite polynomial.\\
$x$, $h$ & $x=k-c_n$ in Section~\ref{tfs:sec:allorders}; $x=\sqrt{L/b}$ and $h$ in Section~\ref{tfs:sec:inverse}.\\
\bottomrule
\end{tabular}
\end{center}
\end{writenote}

\section{Mean matching and covariance}\label{tfs:sec:tilt}
Tilt a single match by $e^{t\cdot\text{total score}}$. Its normalizer and
draw probability are
\[
A(t)=\frac{2e^{3t}+e^{2t}}3,\qquad q=q(t)=\frac1{1+2e^t}.
\]
The decisive outcomes have probability $(1-q)/2$ each. Under this law,
denoted $\Prob_q$, each team has mean
$\mu(q)=(n-1)(3-q)$. For an interior target integer $k$, set
$q=3-k/(n-1)$; its tilted mean is exactly $k$. The change of measure is
\begin{equation}\label{tfs:eq:tilt}
\Prob(X=k\1)=e^{M\log A(t)-nkt}\Prob_q(X=k\1).
\end{equation}
Here the unadorned $\Prob$ always means the original uniform law $q=1/3$.

Put
\[
v(q)=\frac{9-8q-q^2}{4},\qquad \beta(q)=q(1-q),\qquad
\alpha(q)=2v(q).
\]
The covariance matrix has the form
\begin{equation}\label{tfs:eq:sigma}
\Sigma_{n,q}=\lambda_n(q)I+
                 \frac{\rho_n(q)-\lambda_n(q)}n\1\1^T,
\end{equation}
where its transverse and common eigenvalues are
\begin{equation}\label{tfs:eq:eigen}
\lambda_n(q)=n\alpha(q)-\beta(q),\qquad
\rho_n(q)=(n-1)\beta(q).
\end{equation}
Thus the normalized Gaussian integral is
\begin{equation}\label{tfs:eq:J}
J_n(q)=(2\pi)^{-n/2}\lambda_n(q)^{-(n-1)/2}\rho_n(q)^{-1/2}.
\end{equation}
At $q_0=1/3$,
\begin{equation}\label{tfs:eq:baseeig}
\lambda_n=\frac{28n-2}{9},\qquad \rho_n=\frac{2(n-1)}9,
\qquad \alpha=\frac{28}9,\qquad\beta=\frac29.
\end{equation}
Both eigenvalues have order $n$. Using the full score vector at a fixed
common score therefore avoids a singular covariance.

\section{Sharp global localization}\label{tfs:sec:localize}
Let $Q_0$ be a sufficiently small fixed closed interval about $1/3$.
All constants in this section are uniform for $q\in Q_0$. Denote by
$\chi_{n,q}(\theta)$ the characteristic function of the centered total
score vector. At a mean-matching integer score, Fourier inversion gives
\[
\Prob_q(X=k\1)=(2\pi)^{-n}\int_{[-\pi,\pi]^n}\chi_{n,q}(\theta)\,d\theta.
\]

\begin{lemma}[Relative localization]\label{tfs:lem:localize}
There is $c>0$ such that
\begin{equation}\label{tfs:eq:integralbound}
(2\pi)^{-n}\int_{[-\pi,\pi]^n}|\chi_{n,q}(\theta)|\,d\theta
       \leq (1+O(e^{-cn}))J_n(q).
\end{equation}
For every fixed $A>0$, the part of this integral outside
$\max_i|\theta_i|\leq n^{-1/2}\log n$ is $o(n^{-A})J_n(q)$.
\end{lemma}

\subsection{Coercivity on the torus}
The one-match kernel is
\[
\phi_q(x,y)=\frac{1-q}{2}e^{3\ii x}+q e^{\ii(x+y)}
                              +\frac{1-q}{2}e^{3\ii y}.
\]
Let $d(u)=\min_{m\in\Z}|u-2\pi m|$. The three probabilities are bounded
below on $Q_0$. Expanding the squared modulus and using
$1-\cos u\geq c\,d(u)^2$ gives
\begin{equation}\label{tfs:eq:edgecoercive}
|\phi_q(x,y)|^2\leq
\exp\{-c[d(3x-3y)^2+d(2x-y)^2+d(x-2y)^2]\}.
\end{equation}
For distinct $i,j,k$, the identity
\[
\theta_j=(2\theta_i-\theta_j)+(2\theta_j-\theta_k)
                       -(2\theta_i-\theta_k)\pmod{2\pi}
\]
and the squared triangle inequality imply, after summation over ordered
triples,
\[
\sum_{i\ne j}d(2\theta_i-\theta_j)^2
             \geq\frac{n-1}{9}\sum_jd(\theta_j)^2.
\]
Indeed each of the three edge terms is counted $n-2$ times, and each
left-side coordinate is counted $(n-1)(n-2)$ times before cancellation.
Since the double-match modulus is $\prod_{i<j}|\phi_q(\theta_i,\theta_j)|^2$,
\begin{equation}\label{tfs:eq:coercive}
|\chi_{n,q}(\theta)|\leq e^{-cn\sum_i\theta_i^2}
                 \quad(\theta\in[-\pi,\pi]^n,\ n\geq3).
\end{equation}
In particular the origin is the only modulus-one saddle. This coarse
bound alone would lose a factor exponential in $n$ relative to $J_n$;
the next two steps are essential.

\subsection{A negative quartic and an exact Gaussian majorant}
Set $s=(x+y)/2$ and $d=(x-y)/2$. The centered one-match phase is $sU+dV$,
where $(U,V)$ takes the values $(q,3),(q,-3),(q-1,0)$ with probabilities
$(1-q)/2,(1-q)/2,q$. Its fourth cumulant is
\begin{align}\label{tfs:eq:quartic}
\kappa_4(sU+dV)={}&q(1-q)(1-6q+6q^2)s^4\notag\\
 &+54q(1-q)(2q-1)s^2d^2+81(1-q)(3q-2)d^4.
\end{align}
At $q=1/3$ this equals $-2s^4/27-4s^2d^2-54d^4$.
All three coefficients stay strictly negative on a sufficiently small
$Q_0$. Uniform analyticity near the origin and evenness under
$(s,d)\mapsto(-s,-d)$ give
\[
\log|\phi_q|=-\frac{\kappa_2}{2}+\frac{\kappa_4}{24}
                       +O((s^2+d^2)^3).
\]
Choose a fixed $\eta>0$ so that the negative quartic dominates the remainder
when $|x|,|y|\leq\eta$. Multiplying over all matches yields the sharp bound
\begin{equation}\label{tfs:eq:sharp}
|\chi_{n,q}(\theta)|\leq
     e^{-\theta^T\Sigma_{n,q}\theta/2}\qquad
                          (\max_i|\theta_i|\leq\eta).
\end{equation}
Its normalized integral is at most $J_n(q)$, with no dimension-dependent
loss.

\subsection{Outliers outside the fixed box}
Fix $\varepsilon\in(0,1/4)$. If at least $\varepsilon n$ coordinates have
$|\theta_i|>\eta/10$, \eqref{tfs:eq:coercive} gives $e^{-cn^2}$ for the
integrand. Its full contribution is negligible compared with
$J_n(q)=\exp\{-\tfrac n2\log n+O(n)\}$.

Otherwise at least $(1-\varepsilon)n$ coordinates are tiny, meaning
$|\theta_i|\leq\eta/10$. Let $B=\{i:|\theta_i|>\eta\}$, with
$m=|B|\geq1$, and $K=B^c$, with $k=n-m\geq(1-\varepsilon)n$.
For a bad coordinate $x$ and a tiny coordinate $y$,
\[
d(x-2y)\geq d(x)-2|y|\geq\frac45|x|.
\]
Their distinct edges give a factor $e^{-cn\sum_{i\in B}\theta_i^2}$.
Retain every $K$--$K$ edge and apply \eqref{tfs:eq:sharp} to the complete
$k$-team system; bound every remaining factor by $1$. For a fixed set $B$,
the unnormalized integral is consequently at most
\begin{equation}\label{tfs:eq:outlier}
I_k(q)\left[C n^{-1/2}e^{-c'n\eta^2}\right]^m,
\qquad I_k(q)=(2\pi)^kJ_k(q).
\end{equation}
This estimate is valid on the region with enough tiny coordinates; after
the pointwise estimate that restriction can be removed for integration.
There is no extra sum over choices of tiny coordinates.

For completeness, \eqref{tfs:eq:eigen} gives the uniform comparison
\[
I_k(q)\asymp [2\pi/\alpha(q)]^{k/2}k^{-k/2}
\]
for all large $k$ and all $q\in Q_0$.
The residual factors are a bounded multiple of
$(1-\beta/(\alpha k))^{-(k-1)/2}\sqrt{k/(k-1)}$.
It follows that
\[
\frac{I_k(q)n^{-m/2}}{I_n(q)}
 \leq C D^m(n/k)^{k/2}\leq C(De^{1/2})^m.
\]
Summing \eqref{tfs:eq:outlier} over all bad sets gives a relative bound
\[
C\sum_{m\geq1}\binom nm D_1^m e^{-c'nm}=O(e^{-c''n}).
\]
This proves exponentially small relative mass outside the fixed box.

\subsection{Shrinking to the Gaussian scale}
In the fixed box use \eqref{tfs:eq:sharp}. Every coordinate of
$N(0,\Sigma_{n,q}^{-1})$ has variance $O(1/n)$. A union bound gives
\[
\Prob\{\max_i|\theta_i|>n^{-1/2}\log n\}
                    \leq 2n e^{-c\log^2 n}=o(n^{-A})
\]
for every fixed $A$. Together with the outlier bound this proves
Lemma~\ref{tfs:lem:localize}.

\section{The local integral and the score shift}\label{tfs:sec:local}
We first work at $q=1/3$. Under the Gaussian measure associated with
$J_n$, decompose $\theta_i=y_i+a$, where $\sum_i y_i=0$. The vector $y$
and scalar $a$ are independent, and may be represented as
\[
y_i=\frac{Z_i-\overline Z}{\sqrt{\lambda_n}},\qquad
 a=\frac{W}{\sqrt{n\rho_n}},\qquad c=na,
\]
with independent standard normals $Z_1,\ldots,Z_n,W$. Put
$Y_j=\sum_i y_i^j$. The cubic part of the double-match logarithm is exactly
\begin{equation}\label{tfs:eq:cubic}
-\ii\left[\frac{20n+1}{81}Y_3+
           \frac{13n+1}{27}aY_2-\frac{n(n-1)}{81}a^3\right].
\end{equation}
The quartic part at $a=0$ is
\begin{equation}\label{tfs:eq:quartagg}
-\frac{98n-1}{324}Y_4-\frac{89}{108}Y_2^2.
\end{equation}
The quartic terms involving $a$ tend to zero in probability.
On the shrinking box, the sum of degree at least five is
$O(n^2(n^{-1/2}\log n)^5)=o(1)$.

The law of large numbers and central limit theorem give jointly
\begin{align}\label{tfs:eq:limits}
Y_2&\longrightarrow\frac1\alpha,&
 nY_4&\longrightarrow\frac3{\alpha^2},&
 nY_3&\ \Rightarrow\ N(0,6/\alpha^3),&
 c&\ \Rightarrow\ N(0,1/\beta).
\end{align}
The last variable is independent of the others. To verify the cubic limit,
write $H_3(z)=z^3-3z$ and use
\[
\sum_i(Z_i-\overline Z)^3
 =\sum_iH_3(Z_i)-3\overline Z\left(\sum_iZ_i^2-n\right)
                                      +2n\overline Z^3.
\]
After division by $\sqrt n$, the last two terms vanish in probability,
whereas $\Var(H_3(Z))=6$.

The quartic exponent converges to $-561/3136$. The transverse cubic has
limiting variance $25/2058$, and the common cubic converges to
$-\ii(13/84)c$. Their Gaussian expectations therefore give the correction
\begin{equation}\label{tfs:eq:K}
K_0=\exp\left(-\frac{561}{3136}-\frac{25}{4116}
                              -\frac{169}{3136}\right)
    =e^{-7865/32928}.
\end{equation}
Passing to expectations is justified rather than formal: by
\eqref{tfs:eq:sharp}, the complete correction relative to the Gaussian density
has modulus at most $1$ in the fixed box; its shrinking-box complement
has superpolynomially small Gaussian measure. Bounded convergence in
distribution applies to the displayed limiting correction.
The same reasoning holds for every sequence $q\to1/3$.

Now set $\delta=k-\mu_n$. For bounded $\delta$,
\[
q=\frac13-\frac{\delta}{n-1},\qquad
 t=\frac92\frac{\delta}{n-1}+O(n^{-2}),
\]
and Taylor expansion gives
\begin{align}
M\log A(t)-nkt&=-\frac94\delta^2+o(1),\label{tfs:eq:rateleading}\\
\frac{J_n(q)}{J_n(1/3)}&=\exp\left(-\frac{39}{56}\delta+o(1)\right).
                                                       \label{tfs:eq:detleading}
\end{align}
For the determinant, note
$\lambda'_n(1/3)=-(13n+1)/3$. Its $n-1$ transverse factors produce the
finite linear term; the common eigenvalue changes only by relative
$O(1/n)$. Combining \eqref{tfs:eq:K}--\eqref{tfs:eq:detleading} and completing the
square proves
\begin{equation}\label{tfs:eq:fixedscore}
\Prob(X=k\1)=(1+o(1))J_n(1/3)e^{-12181/65856}
                e^{-9(\delta+13/84)^2/4}.
\end{equation}
The error is uniform for bounded $\delta$: otherwise a violating sequence
has a convergent subsequence of $\delta$ and the preceding sequential
argument gives a contradiction.

\section{Dominated summation and the conditional common score}\label{tfs:sec:summation}
Pointwise fixed-score equivalents cannot be summed without a uniform
tail estimate. If $|\delta|\leq\varepsilon_0 n$, with $\varepsilon_0$
sufficiently small, the exact mean-matching tilt lies in $Q_0$.
Lemma~\ref{tfs:lem:localize} implies $\Prob_q(X=k\1)\leq C J_n(q)$.
The exact rate and determinant satisfy
\[
M\log A(t)-nkt\leq-c\delta^2,\qquad
\frac{J_n(q)}{J_n(1/3)}\leq C e^{C|\delta|}.
\]
For the first bound, $(\log A)''=q(1-q)$ is bounded above and below on
the tilt interval; the second follows from the explicit eigenvalues.
Thus
\begin{equation}\label{tfs:eq:majorant}
\frac{\Prob(X=k\1)}{J_n(1/3)}\leq
                         C e^{-c\delta^2+C|\delta|}.
\end{equation}
If $|\delta|>\varepsilon_0 n$, the tied event forces the total score
$nk$ to differ from its mean by at least $\varepsilon_0 n^2$.
The number of draws is $\mathrm{Binomial}(M,1/3)$ and the total score is
$3M$ minus this number. A binomial Chernoff bound makes the sum of these
terms at most $O(n)e^{-cn^2}=o(J_n(1/3))$.

Dominated summation on each residue subsequence now gives
\begin{equation}\label{tfs:eq:sumresult}
\Prob(\text{all tied})\sim J_n(1/3)e^{-12181/65856}
                              \sum_{k\in\Z}e^{-9(k-c_n)^2/4}.
\end{equation}
The natural endpoint scores are in the negligible tail, so replacing the
finite score range by $\Z$ loses no mass at this scale. Equations
\eqref{tfs:eq:poisson}, \eqref{tfs:eq:baseeig}, and \eqref{tfs:eq:sumresult} give
\eqref{tfs:eq:exactcarrier}. Finally,
\[
\left(\frac{28n-2}{28n}\right)^{-(n-1)/2}\longrightarrow e^{1/28},
\qquad \sqrt{\frac n{n-1}}\longrightarrow1,
\]
and $1/28-12181/65856=-9829/65856$, proving
Theorem~\ref{tfs:thm:leading}.

\begin{proposition}[Conditional score law]\label{tfs:prop:conditional}
On each residue subsequence, after an integer translation, the common
score conditioned on all teams being tied converges to the discrete
Gaussian with weights proportional to $e^{-9(k-c_n)^2/4}$. Every fixed
moment converges as well. In particular,
\begin{align}\label{tfs:eq:condmean}
\E[k\mid\text{all tied}]-\mu_n
 &=-\frac{13}{84}+\frac29\frac{\Theta'(c_n)}{\Theta(c_n)}+o(1),\\
\Var(k\mid\text{all tied})
 &=\frac29+\frac4{81}(\log\Theta)''(c_n)+o(1).\label{tfs:eq:condvar}
\end{align}
\end{proposition}
\begin{proof}
The same majorant \eqref{tfs:eq:majorant} permits multiplication by every
fixed power of $|\delta|$. If $Z(c)=\sum_k e^{-9(k-c)^2/4}$, differentiation
gives $\E(k-c)=\tfrac29(\log Z)'$ and
$\Var(k)=\tfrac29+\tfrac4{81}(\log Z)''$. Use
$Z(c)=2\sqrt\pi\Theta(c)/3$.
\end{proof}
The parameter $c_n$ is a Gaussian center, generally not the actual
conditional mean. The shift $-13/84$ comes from the mixed common--transverse
cubic; omitting it changes the leading theta multiplier. On the tied event,
the number of draws is exactly $n[3(n-1)-k]$ and hence is a multiple of $n$.

\section{Expansion to every fixed order}\label{tfs:sec:allorders}
We now establish a uniform local expansion for $q\in Q_0$, then perform
the tilted score sum. The external input is Isaev's complex cumulant
theorem \cite[Theorem 1.2]{isaev}. We state the form used here to make the
hypothesis verification explicit.

\subsection{The complex cumulant input}
Let $Z_1,\ldots,Z_n$ be independent on a product domain $\Omega$ and let
$f:\Omega\to\mathbb C$ be bounded. For coordinate replacement
$z\triangleright_W y$, define
\[
\overline\Delta_V(f)=\sup_{y\in\Omega}
 \left|\E\sum_{W\subseteq V}(-1)^{|W|}f(Z\triangleright_Wy)\right|,
\qquad S_k(f)=\max_j\sum_{\substack{|V|=k\\j\in V}}
                                            \overline\Delta_V(f).
\]
For a fixed integer $m\geq1$, suppose $0\leq\gamma\leq1/100$ and
\begin{equation}\label{tfs:eq:isaevhyp}
\sum_{k=t}^n e^{3\gamma k/2}2^t\binom{k-1}{t-1}S_k(f)\leq\gamma
                  \qquad(1\leq t\leq m).
\end{equation}
Then
\begin{equation}\label{tfs:eq:isaevconcl}
\E e^{f(Z)}=(1+\zeta)^n
       \exp\left(\sum_{j=1}^m\frac{\kappa_j(f(Z))}{j!}\right),
\qquad |\zeta|\leq e^{(100\gamma)^{m+1}}-1.
\end{equation}
For $2\leq j\leq m$ the theorem also gives
\[
|\kappa_j(f(Z))|\leq\frac{n(j-1)!(80\gamma)^j}{50j}.
\]
These are algebraic complex cumulants without conjugation.
Bounding each $\overline\Delta_V$ by its supremum mixed difference is
sufficient for \eqref{tfs:eq:isaevhyp}.

\subsection{Symmetric whitening and independent truncation}
The symmetric inverse square root of \eqref{tfs:eq:sigma} is
\begin{equation}\label{tfs:eq:T}
T=\lambda^{-1/2}I+
           \frac{\rho^{-1/2}-\lambda^{-1/2}}n\1\1^T,
\qquad\theta=TZ,
\end{equation}
with independent standard normals $Z_i$. Diagonal entries of $T$ are
$O(n^{-1/2})$, off-diagonal entries are $O(n^{-3/2})$, and
\begin{equation}\label{tfs:eq:rows}
\max_i\sum_j|T_{ij}|+\max_j\sum_i|T_{ij}|\leq Cn^{-1/2}
\end{equation}
uniformly in $q$. This spreads the common fluctuation among the $n$
independent variables; it does not introduce a single coordinate with
order-one influence in the cubic exponent.

Take $K=\log n$ and condition each $Z_i$ on $|Z_i|\leq K$. Independence is
preserved. The omitted probability is smaller than every inverse power
of $n$, and \eqref{tfs:eq:rows} puts the image of the product cube inside
$\max_i|\theta_i|\leq CK/\sqrt n<\eta$. Conversely, removing from the
fixed $\theta$-box the part outside this image costs at most the omitted
Gaussian probability, by \eqref{tfs:eq:sharp}. The normalization introduced
by conditioning also differs from $1$ by a superpolynomially small amount.

Let $L_r(\theta)$ be the degree-$r$ term in the logarithm of the centered
characteristic function, with the quadratic removed, and put
\[
f_D(Z)=\sum_{r=3}^D L_r(TZ).
\]
Each $L_r$ is a sum over edges of bounded-coefficient homogeneous
bivariate polynomials. On the product cube its Taylor error is
\begin{equation}\label{tfs:eq:taylorerror}
O_D\bigl(n^2(K/\sqrt n)^{D+1}\bigr)
                 =O_D(n^{(3-D)/2}K^{D+1}).
\end{equation}
The exact correction has modulus at most $1$, and the truncated correction
has modulus at most the exponential of this error. Hence replacement of
the exponent changes the normalized integral by the same order.

\subsection{Mixed derivative bounds}
For $1\leq k\leq r$ and a fixed index $j_1$, the chain rule gives
\begin{align}\label{tfs:eq:derivative}
&\sum_{j_2,\ldots,j_k}\sup_{|Z_i|\leq K}
 |\partial_{j_1}\cdots\partial_{j_k}L_r(TZ)|\notag\\
&\quad\leq C_r(K/\sqrt n)^{r-k}
 \left[\sum_{uv\text{ edge}}(|T_{u,j_1}|+|T_{v,j_1}|)\right]
 \left[\max_{uv\text{ edge}}\sum_j(|T_{u,j}|+|T_{v,j}|)\right]^{k-1}
 \notag\\
&\quad\leq C_r n^{1-r/2}K^{r-k}.
\end{align}
The first bracket is $(n-1)$ times a column absolute sum, hence
$O(\sqrt n)$; the second is $O(n^{-1/2})$. Including repeated indices
only enlarges the bound, so it controls the subset sums required above.
A $k$-fold mixed difference on the cube is bounded by $(2K)^k$ times its
mixed derivative. Consequently
\begin{equation}\label{tfs:eq:Sk}
S_k(f_D)\leq C_D n^{-1/2}K^D\quad(1\leq k\leq D),\qquad
S_k(f_D)=0\quad(k>D).
\end{equation}
For fixed $D,m$, choose
$\gamma=C_{D,m}n^{-1/2}K^D$ with a sufficiently large constant. There
are only finitely many nonzero terms in \eqref{tfs:eq:isaevhyp}; it holds for
all large $n$, as does $\gamma\leq1/100$. The relative error in
\eqref{tfs:eq:isaevconcl} is therefore
\begin{equation}\label{tfs:eq:cumulanttail}
O_{D,m}(n\gamma^{m+1})=
                O_{D,m}(n^{(1-m)/2}K^{D(m+1)}).
\end{equation}
Choosing fixed $D,m$ sufficiently large gives any requested inverse-power
precision. In particular, this resums the order-one cubic fluctuation;
no finite Taylor expansion of its exponential is being presumed valid.

Truncated-normal cumulants can be replaced by full-normal cumulants with
superpolynomial error. Every fixed moment of $f_D$ is a fixed-degree
Gaussian polynomial with coefficients growing at most polynomially in
$n$. Cauchy--Schwarz and Gaussian polynomial moment bounds make its tail
contribution at most $n^C e^{-c\log^2 n}$. Cumulants are finite polynomial
combinations of moments. Thus, to any prescribed inverse-power precision,
the local correction equals
\begin{equation}\label{tfs:eq:localexpansion}
\exp\left(\sum_{j=1}^m\frac{\kappa_j(f_D(Z))}{j!}\right)
                    (1+\text{controlled error}),
\end{equation}
where now $Z$ is an unconditioned standard Gaussian vector. The leading
limit is positive uniformly on a sufficiently small $Q_0$, so the
additive Taylor and localization errors can also be expressed relatively.

\subsection{Exact finite coefficients and integer powers}
For clarity use $u,v$ for the two covariance entries:
\begin{equation}\label{tfs:eq:coventries}
u=\lambda^{-1},\qquad
v=\frac{\rho^{-1}-\lambda^{-1}}n,\qquad
\E\theta_i\theta_j=u\,\1_{i=j}+v.
\end{equation}
They are rational functions of $n,q$. For fixed degree, every moment
of products of the $L_r$ is rational in $n,q$ by Wick's formula: grouping
vertex identifications by a finite set partition yields a falling
factorial $(n)_\ell$ times a finite Gaussian pairing sum.
Here is a completely explicit rule. Let $P_r=\sum_i\theta_i^r$, $P_0=n$.
For a product $P_{r_1}\cdots P_{r_t}$, sum over all set partitions
$\pi$ of $\{1,\ldots,t\}$. Put $d_B=\sum_{j\in B}r_j$ for each block
$B$, $d=\sum_B d_B$, and $h=\sum_B h_B$. If $d$ is odd the moment is zero.
If $d$ is even the contribution of $\pi$ is
\begin{equation}\label{tfs:eq:wick}
(n)_{|\pi|}
\sum_{0\leq h_B\leq\lfloor d_B/2\rfloor}
\left[\prod_{B\in\pi}\binom{d_B}{2h_B}(2h_B-1)!!\right]
(d-2h-1)!!\,u^h v^{d/2-h}.
\end{equation}
Here $(-1)!!=1$. To prove the rule, for distinct block indices represent
$\theta_B=\sqrt u\,Z_B+\sqrt v\,W$ and expand, using Gaussian even moments.
On $Q_0$, $v>0$; the resulting polynomial identity is algebraic regardless
of this representation. Factors $P_0$ can simply be taken out as powers
of $n$. Every symmetric bivariate edge polynomial is a finite combination
of the $P_r$, since
\[
\sum_{i\ne j}\theta_i^a\theta_j^b=P_aP_b-P_{a+b}.
\]

If $M_j=\E f_D^j$, with $M_0=1$, the ordinary cumulant recursion is
\begin{equation}\label{tfs:eq:cumrec}
\kappa_j=M_j-\sum_{i=1}^{j-1}\binom{j-1}{i-1}\kappa_i M_{j-i}.
\end{equation}
Equations \eqref{tfs:eq:wick}--\eqref{tfs:eq:cumrec} give a finite rational
algorithm. The first cumulant is $O(1)$: odd-degree means vanish and
even $r\geq4$ have mean $O(n^{2-r/2})$. Isaev's bound controls higher
cumulants by $O(n^{1-j/2}\log^C n)$. Since they are rational functions
with integer-power Laurent expansions at infinity, no positive powers
can occur, and no constant term can occur for $j\geq3$.
Thus the logarithm in \eqref{tfs:eq:localexpansion} has an ordinary $1/n$
expansion with coefficients regular uniformly on $Q_0$.

For a requested grade $J$, a finite procedure is:
\begin{enumerate}
\item Choose $D\geq2J+8$ and $m\geq2J+5$. Both
\eqref{tfs:eq:taylorerror} and \eqref{tfs:eq:cumulanttail} are
$o(n^{-J-1})$.
\item Compute the one-match cumulants through degree $D$ from the
three-point distribution, and sum the double-match logarithmic terms.
\item Use \eqref{tfs:eq:wick} and \eqref{tfs:eq:cumrec} to form
$\sum_{j\leq m}\kappa_j(f_D)/j!$ exactly.
\item Expand this rational function through $n^{-J-1}$ and exponentiate
as a finite formal series. Retaining grades through $J$ leaves
$O_J(n^{-J-1})$.
\end{enumerate}
Increasing $D,m$ cannot change the resulting coefficients, by uniqueness
of asymptotic power series for the same local integral.

\subsection{From local series to theta derivatives}
Substitute $q=1/3-\delta/(n-1)$ into the uniform expansion and multiply
by the exact tilt factor and determinant. Each coefficient is a
polynomial in $\delta$ times
$e^{-9(\delta+13/84)^2/4}$. On $|\delta|\leq\log n$, Taylor remainders
are bounded by $n^{-J-1}$ times a fixed polynomial in $|\delta|$ times
$e^{-c\delta^2+C|\delta|}$. The sum of this bound is uniform; the part
$|\delta|>\log n$ is superpolynomially negligible by
\eqref{tfs:eq:majorant} and the large-deviation bound. This proves the
remainder in \eqref{tfs:eq:allorders}.

Finally put $x=k-c_n$, so $\delta=x-13/84$.
Differentiating $Z(c)=\sum_k e^{-9(k-c)^2/4}$ expresses every Gaussian
power sum as a rational linear combination of derivatives of $Z$:
for example, $Z'=\tfrac92\sum_k(k-c)e^{-9(k-c)^2/4}$.
Recursing through higher derivatives yields the claimed structure of
$H_j$ in Theorem~\ref{tfs:thm:allorders}. On the three residue subsequences,
these functions take constant values.

\section{The first correction}\label{tfs:sec:first}
The preceding argument proves existence and error control at every fixed
order. We now give the finite algebra for its first nonconstant term.
At $q=1/3$, the necessary logarithmic polynomials are
\begin{align}
L_3={}&\ii\left(\frac7{27}P_1P_2-\frac{20n+1}{81}P_3\right),\label{tfs:eq:L3}\\
L_4={}&\frac{91}{81}P_1P_3-\frac{89}{108}P_2^2
                        -\frac{98n-1}{324}P_4,\label{tfs:eq:L4}\\
L_5={}&\ii\left(\frac{407}{972}P_1P_4-\frac{133}{486}P_2P_3
                        -\frac{140n+1}{972}P_5\right),\label{tfs:eq:L5}\\
L_6={}&\frac{7861}{14580}P_1P_5-\frac{8015}{5832}P_2P_4
       +\frac{8077}{8748}P_3^2-\frac{7718n-7}{87480}P_6.\label{tfs:eq:L6}
\end{align}
Take $D=6,m=4$. The Taylor and cumulant-tail errors are
$O(n^{-3/2}\log^C n)=o(1/n)$. Through grade $1/n$, the logarithmic
correction consists of
\begin{align}\label{tfs:eq:seven}
&\E L_4+\tfrac12\E L_3^2+\E L_6+\E(L_3L_5)
 +\tfrac12\Var(L_4)\notag\\
&\hspace{35mm}+\tfrac12\kappa(L_3,L_3,L_4)
                         +\tfrac1{24}\kappa_4(L_3).
\end{align}
All variances and cumulants here are algebraic, without complex
conjugation.

\subsection{Why the omitted mixtures are smaller}
In every fixed Gaussian $L^p$ norm,
\[
L_3=O(1),\quad L_4-\E L_4=O(n^{-1/2}),\quad
L_5=O(n^{-1}),\quad L_6-\E L_6=O(n^{-3/2}).
\]
These assertions follow by inserting $\theta_i=y_i+c/n$ into
\eqref{tfs:eq:L3}--\eqref{tfs:eq:L6} and using Gaussian polynomial moment bounds.
More specifically,
\[
\bigl(L_3,\ \sqrt n(L_4-\E L_4),\ nL_5\bigr)
\]
converges jointly, with all fixed moments, to a centered Gaussian vector.
The first and third coordinates are linear combinations of the common
normal $c$ and empirical Hermite sums of degrees $3$ and $5$; the second
is a linear combination of Hermite sums of degrees $2$ and $4$.
The even component is independent of the odd components in this
Gaussian limit. Also $nL_6$ tends to a deterministic constant.

Consequently the potentially order-$1/n$ mixtures
$\kappa(L_3,L_3,L_4,L_4)$ and
$\kappa(L_3,L_3,L_3,L_5)$ are actually $o(1/n)$: their scaled limits are
fourth cumulants of a Gaussian vector. Parity kills the odd total-degree
terms exactly. The displayed $L^p$ estimates control the other omissions;
cumulants of order above four are covered by \eqref{tfs:eq:cumulanttail}.
This establishes \eqref{tfs:eq:seven} without supposing $L_3=o(1)$.

\subsection{Exact local coefficient}
The finite Wick calculation gives the following constants and $1/n$
coefficients. The entries are exact rational numbers.
\begin{center}
\begin{tabular}{@{}lrr@{}}
\toprule
Quantity & Constant & Coefficient of $1/n$\\
\midrule
$\E L_4$ & $-561/3136$ & $159/5488$\\
$\E L_6$ & $0$ & $-31763/175616$\\
$\E L_3^2$ & $-3949/32928$ & $14407/115248$\\
$\E(L_3L_5)$ & $0$ & $-26561/175616$\\
$\Var(L_4)$ & $0$ & $343533/1229312$\\
$\kappa(L_3,L_3,L_4)$ & $0$ & $492179/2151296$\\
$\kappa_4(L_3)$ & $0$ & $4310925/15059072$\\
\bottomrule
\end{tabular}
\end{center}
Thus
\begin{equation}\label{tfs:eq:localfirst}
\log(\text{local correction at }q=1/3)
=-\frac{7865}{32928}+\frac{1147975}{45177216}\frac1n+O(n^{-2}).
\end{equation}
The initial calculation gives $o(1/n)$; the already proved all-order
expansion upgrades this to $O(n^{-2})$. A second raw-moment aggregation
through cumulant order six yields the same two coefficients and zero
constant and $1/n$ terms for the fifth and sixth cumulants.
A further Gaussian integration-by-parts computation uses the coordinate
covariance directly and reproduces these first-correction moments and the
vanishing coefficients of the two omitted fourth-cumulant mixtures.

\subsection{Tilt and carrier contributions}
For general $q$, put
\[
d_3=\frac{\beta(q)(q+13)}{12},\qquad
b_3=\frac{\beta(q)(q+4)}2,
\]
and let $A_4,B_4,C_4$ be the three quartic coefficients in
\eqref{tfs:eq:quartic}. The leading local logarithmic correction is
\begin{align}\label{tfs:eq:Cq}
C(q)&=\frac{3A_4+B_4+3C_4}{96\alpha(q)^2}
       -\frac{3d_3^2}{\alpha(q)^3}
       -\frac{b_3^2}{2\beta(q)\alpha(q)^2}\notag\\
&=-\frac{q^4-28q^3-122q^2+1944q-2187}
                    {12(q-1)(q+9)^3}.
\end{align}
In particular $C(1/3)=-7865/32928$ and $C'(1/3)=531/153664$.
The dependence of the leading local correction on the tilt contributes
$-(531/153664)\delta/n$.

The exact tilt rate is $-n(n-1)D(q\Vert1/3)$, where $D$ is Bernoulli
relative entropy. Its expansion and that of the determinant are
\begin{align}
-n(n-1)D(q\Vert1/3)
 &=-\frac94\delta^2+\frac1n\left(-\frac94\delta^2-\frac98\delta^3\right)
                         +O(n^{-2}\poly(\delta)),\label{tfs:eq:ratefirst}\\
\log\frac{J_n(q)}{J_n(1/3)}
 &=-\frac{39}{56}\delta+
      \frac1n\left(\frac{507}{784}\delta+\frac{1773}{3136}\delta^2\right)
                         +O(n^{-2}\poly(\delta)).\label{tfs:eq:detfirst}
\end{align}
Combining these with \eqref{tfs:eq:localfirst} gives, relative to the exact
$J_n(1/3)$ shifted-Gaussian carrier, the first correction
\[
\frac{1147975}{45177216}
+\frac{98841}{153664}\delta-\frac{5283}{3136}\delta^2-\frac98\delta^3.
\]
The logarithm of the remaining covariance normalization factor is
\[
\log\left[\sqrt{\frac n{n-1}}
             \left(1-\frac1{14n}\right)^{-(n-1)/2}\right]
           =\frac1{28}+\frac{365}{784n}+O(n^{-2}).
\]
The constant $1/28$ is already included in $B_n$.
Adding $365/784$ to the displayed polynomial's constant gives
\eqref{tfs:eq:Q}. Dominated summation proves \eqref{tfs:eq:Rr}--\eqref{tfs:eq:first}.

\section{Residue-aware asymptotic inversion}\label{tfs:sec:inverse}
Let $b=\log3$ and set
\begin{equation}\label{tfs:eq:de}
d=-b+\tfrac12\log\frac9{56\pi},\qquad
e_r=-\tfrac12\log\frac9{56\pi}-\frac{9829}{65856}+\log\Theta_r,
\end{equation}
where $\Theta_r=\Theta(c_n)$ for $n\equiv r\pmod3$.
Taking logarithms of Theorem~\ref{tfs:thm:allorders},
\begin{equation}\label{tfs:eq:logforward}
\log a(n)=bn^2-\tfrac n2\log n+dn+e_r
  +\sum_{j=1}^R\ell_{j,r}n^{-j}+O_R(n^{-R-1}),\qquad
\ell_{1,r}=R_r.
\end{equation}
There is no standalone $\log n$ term: the two logarithms from
$n^{-1/2}$ and $n^{-(n-1)/2}$ cancel it exactly.

For a fixed $R$, define the explicit smooth function
\begin{equation}\label{tfs:eq:F}
F_{R,r}(t)=bt^2-\tfrac t2\log t+dt+e_r
                         +\sum_{j=1}^R\ell_{j,r}t^{-j}.
\end{equation}
Its derivative is $2bt+O(\log t)>0$ for large $t$, so its large real inverse
$\nu_{R,r}(L)$ exists uniquely. This differentiates the explicit
approximant, not an unspecified asymptotic remainder.
For example, put
\[
x=\sqrt{L/b},\qquad h=\frac{\log x}{4b}-\frac d{2b}.
\]
The first three terms of the smooth inverse are
\begin{equation}\label{tfs:eq:inverse}
\nu_{0,r}(L)=x+h+
 \frac{-bh^2+\tfrac h2(\log x+1)-dh-e_r}{2bx}+o(x^{-1}).
\end{equation}
Successive substitution or Newton coefficient matching produces any
fixed further grade in $x^{-1}$, with polynomial coefficients in
$\log x$. The derivative bound for $F_{R,r}$ controls the remainder.

\begin{proposition}[Integer threshold enclosure]\label{tfs:prop:inverse}
For fixed $R$ and all sufficiently large $L$, let
$\nu_r=\nu_{R,r}(L)$ and $\epsilon_r=C\nu_r^{-R-2}$, with $C$
sufficiently large. Define
\[
N_r^- =r+3\ceil{\frac{\nu_r-\epsilon_r-r}{3}},\qquad
N_r^+ =r+3\ceil{\frac{\nu_r+\epsilon_r-r}{3}}.
\]
If $N(L)=\min\{n\geq1:a(n)\geq e^L\}$, then
\begin{equation}\label{tfs:eq:enclose}
\min_{r=0,1,2}N_r^-\ \leq\ N(L)\ \leq\ \min_{r=0,1,2}N_r^+.
\end{equation}
If a truncated inverse series replaces $\nu_r$, its independently bounded
absolute error must be added to $\epsilon_r$.
\end{proposition}
\begin{proof}
By \eqref{tfs:eq:logforward}, the forward error near $\nu_r$ is
$O(\nu_r^{-R-1})$. The derivative of \eqref{tfs:eq:F} is bounded below by a
positive multiple of $\nu_r$ there. Hence values below
$\nu_r-C\nu_r^{-R-2}$ cannot cross the threshold, and values above
$\nu_r+C\nu_r^{-R-2}$ do cross it, with $C$ large enough.
Taking the next integer in residue class $r$ proves the two bounds for
that residue and then \eqref{tfs:eq:enclose}. Eventual strict monotonicity
also follows directly from the integer forward formula: the successive
main-term difference is $2bn+O(\log n)$, the residue constants are bounded,
and the errors tend to zero. Finite initial exceptions do not affect
large thresholds.
\end{proof}
The asymptotic $O$-constants and their onset are not numerically certified
here. These are eventual theoretical enclosures, not ready-made finite
input certificates. If a branch lies within its uncertainty of the
residue lattice, no floor-free asymptotic formula resolves the rounding.

\begin{writenote}[Added 2 October 2026, batch~77: relation to repository results]
The inversion of this section falls under the general inversion apparatus
of ProveIt's transseries volumes \cite{tfs-tai,tfs-cti}, and no novelty is
claimed for the method. The dominant block here is quadratic, $bt^2$, so
the Lambert
cores (\texttt{p0:thm:lambert-core}, \texttt{p0:prop:factorial-core}) and
the linear--logarithmic reversion \texttt{p0:thm:lambert-centered} of
\cite{tfs-tai} do not apply: the seed is $x=\sqrt{L/b}$ and no Lambert
function occurs. The phase $bt^2-\tfrac t2\log t+dt$ is an admissible core
in the sense of \texttt{p0:def:core} of \cite{tfs-tai} (its slope tends to
$+\infty$), and the successive substitution or Newton coefficient matching
described after \eqref{tfs:eq:inverse} is the formal reversion about such a
core, \texttt{p0:thm:core-reversion} of \cite{tfs-tai} (in \cite{tfs-cti},
the scaled local reversion \texttt{t2:prop:scaled}, stated for an arbitrary
core). The residue-aware enclosure of Proposition~\ref{tfs:prop:inverse}
combines two general parts of \texttt{p0:thm:staircase}: part~(4), the
residue-class formula $\min_\varrho\bigl(\varrho+m\lceil(\nu_\varrho-\varrho)/m\rceil\bigr)$
(written there with modulus $r$ and residue $\rho$; here the modulus is
$m=3$, the residue is called $r$, and $\rho_n$ is an eigenvalue), which is
the shape of $N_r^{\pm}$ and \eqref{tfs:eq:enclose}; and the
separation condition of part~(2), applied in each residue class with the
width $\epsilon_r$, because $F_{R,r}$ is a smooth approximant and not an
admissible interpolation of $a(n)$ on the class. The closing remark above
is the failure mode of part~(2). The arithmetic of those two parts is
formalized in Lean as \texttt{Fabius.staircase\_separation} and
\texttt{Fabius.isLeast\_residue\_class} in
\nolinkurl{Analysis/FabiusFunction/Lean/FabiusFunction/StaircaseInversion.lean};
those statements concern arbitrary real numbers and monotone functions and
give no formal status to anything in this section. Specific to this
report are the forward expansion \eqref{tfs:eq:logforward}, with its three
residue constants $e_r$ and $\ell_{1,r}=R_r$ and the exact cancellation of
the standalone $\log n$ term, and the explicit terms of
\eqref{tfs:eq:inverse}.
\end{writenote}

\section{Prior work and scope}\label{tfs:sec:prior}
The football model and its exact finite enumeration belong to
Jehn, Habermann and Lavrov \cite{football}; the present argument concerns
its large-$n$ asymptotics. A bounded search of the exact sequence, title,
and related football-score and dense-graph asymptotic literature,
completed on 2 October 2026, located no football-specific asymptotic
for this count. That negative search is not an exhaustive novelty
certificate.

Several substantial general precedents are relevant. Barvinok and
Hartigan \cite{barvinok} give maximum-entropy Edgeworth estimates, with
hypotheses that include local correction and global localization.
Isaev and McKay \cite{martingale} treat complex perturbed Gaussian
integrals. Isaev, McKay and Zhang \cite{eulerian} develop cumulant
expansions for Eulerian orientations; their real-valued theorem cannot
simply be applied to the present complex football exponent.
Isaev's complex theorem \cite{isaev} is the explicit higher-order input
used in Section~\ref{tfs:sec:allorders}. The graph-factor work of Isaev and
McKay \cite{factors} provides a close methodological precedent for
whitening and derivative estimates, but its Bernoulli edge model does
not immediately give the three-outcome football count.

Accordingly, no new general local-limit or cumulant theorem is claimed.
The model-specific work consists of the sharp localization, the common
score tilt and lattice summation, their constants and shift, the checked
application of the complex theorem, and the resulting explicit first
correction and inverse formulation.

\section{Reproducibility and limits}\label{tfs:sec:replay}
The accompanying source package contains this editable LaTeX report,
the compiled PDF, portable Python scripts, frozen expected JSON results,
and a SHA256 manifest. From the extracted package root, install the
listed dependencies and run \texttt{python replay.py}. Generated outputs
are written separately from the expected fixtures. The package-local
\texttt{build-report.sh} rebuilds the PDF with a standard TeX installation.
The README gives exact commands and a map of the files.

The replay checks one-match cumulants, symmetric-polynomial identities,
leading constants, exact counts through four teams, the finite Wick
calculation for the first correction, an alternative raw-cumulant
aggregation, a Gaussian integration-by-parts calculation, residue
constants and conditional moments, and selected finite-$n$
Gaussian-cumulant diagnostics.
The known counts for five through eight teams are comparison data from
\cite{football,oeis}, not recomputed exact enumerations in this replay.
The generic every-fixed-order algorithm is specified mathematically in
\eqref{tfs:eq:wick}--\eqref{tfs:eq:cumrec}; the supplied programs implement the
coefficient checks stated above, rather than an optimized arbitrary-order
symbolic engine.

Finite algebraic checks do not replace the localization, error bounds,
or convergence arguments. The proof establishes every fixed order, not
convergence of the infinite asymptotic series or uniformity when the
truncation order grows with $n$. Numerical theta and correction values
are diagnostic decimals; no interval-certified constants, effective
remainder constants, or certified finite-$n$ approximation thresholds
are asserted. Small-$n$ comparisons are not evidence of an error bound.
Only original report and verification material is packaged; cited papers
remain available at their primary source links.

\begin{writenote}[Added 2 October 2026, batch~77]
The source package described above is not shipped as such. In the
collection the nine programs, \texttt{replay.py} and
\texttt{build-\allowbreak report.sh} are in \texttt{code/}, and the seven
frozen fixtures \texttt{expected/\allowbreak<name>.json} and
\texttt{requirements.txt} in \texttt{data/} (as
\texttt{data/\allowbreak expected-\allowbreak<name>.json}), all
byte-identical to the delivery; the compiled PDF and the SHA256 manifest
\texttt{MANIFEST.sha256} (verified at intake and retired) are not shipped
and survive in the archive. \texttt{replay.py} locates the checks under
\texttt{code/} and the fixtures under \texttt{expected/} relative to its
own directory, so in the collection layout it finds neither;
\texttt{build-\allowbreak report.sh} compiles
\texttt{report/\allowbreak tied-\allowbreak football-\allowbreak asymptotics.tex},
which is shipped only as this \texttt{article.tex}. The individual check
scripts write their default output beside the report unless directed
elsewhere. The report's \texttt{README.md} gives a replay on a scratch
copy, rebuilt in the delivered layout from the shipped files, that leaves
them untouched.
\end{writenote}

\begingroup\raggedright
\begin{thebibliography}{9}
\bibitem{football}
R. Jehn, K. Habermann, and M. Lavrov,
\emph{Number of ways that a football league can complete with all teams
having the same number of points}, arXiv:2503.14509v1 (2025).
\url{https://arxiv.org/abs/2503.14509}.

\bibitem{oeis}
OEIS Foundation Inc., \emph{The On-Line Encyclopedia of Integer Sequences},
entry A380592. \url{https://oeis.org/A380592}.

\bibitem{isaev}
M. Isaev, \emph{A tail bound for cumulant series for complex functions of
independent random variables}, arXiv:2508.16952v2 (2025), Theorem 1.2.
\url{https://arxiv.org/abs/2508.16952v2}.

\bibitem{barvinok}
A. Barvinok and J. A. Hartigan, \emph{Maximum entropy Edgeworth estimates
of the number of integer points in polytopes}, arXiv:0910.2497v2,
Theorem 1. \url{https://arxiv.org/abs/0910.2497}.

\bibitem{eulerian}
M. Isaev, B. D. McKay, and R.-R. Zhang,
\emph{Cumulant expansion for counting Eulerian orientations},
Journal of Combinatorial Theory, Series B \textbf{172} (2025), 263--314.
\url{https://arxiv.org/abs/2309.15473v2}.

\bibitem{martingale}
M. Isaev and B. D. McKay, \emph{Complex martingales and asymptotic enumeration},
Random Structures \& Algorithms \textbf{52} (2018), 617--661.
\url{https://arxiv.org/abs/1604.08305}.

\bibitem{factors}
M. Isaev and B. D. McKay,
\emph{Asymptotic enumeration of graph factors by cumulant expansion},
arXiv:2508.18731v1 (2025).
\url{https://arxiv.org/abs/2508.18731}.

\bibitem{tfs-tai} [write] ProveIt, \emph{Transseries and inversion}, the
transseries-and-inversion volume (\texttt{p0:def:core},
\texttt{p0:thm:core-reversion}, \texttt{p0:thm:staircase}; not
applicable: \texttt{p0:thm:lambert-core}, \texttt{p0:prop:factorial-core},
\texttt{p0:thm:lambert-centered}),
\nolinkurl{Analysis/Transseries/docs/series-and-transseries/Transseries_And_Inversion/transseries_and_inversion.tex}.

\bibitem{tfs-cti} [write] ProveIt, \emph{Combinatorial transseries
inverses} (\texttt{t2:prop:scaled}),
\nolinkurl{Analysis/Transseries/docs/series-and-transseries/Combinatorial_Transseries_Inverses/Combinatorial_Transseries_Inverses.tex}.
\end{thebibliography}
\endgroup
\end{document}
